\section{Data Synthesis}
\label{sec:data_synthesis}

The efficacy of the policy learning framework (Section~\ref{sec:agentic_training_framework}) hinges on the availability of high-quality training trajectories. However, training a robust agent for autonomous data preparation presents a significant \textit{data scarcity} challenge.
Existing benchmarks are often limited in scope, focusing on isolated sub-tasks such as simple data cleaning or table reconstruction with short operator sequences. In contrast, real-world data preparation is a composite problem that necessitates the orchestration of cleaning, transformation, and structural reshaping~\cite{heer2015predictive} within a single, long-horizon pipeline~\cite{2011-wrangler}.

To bridge this gap, we propose a scalable data synthesis engine. We leverage the rich semantic diversity of public Text-to-SQL datasets (e.g., Spider~\cite{spider_dataset}, Bird~\cite{bird_dataset}), which accurately reflect real-world analytical intents.
We transform these SQL-oriented benchmarks into autonomous data preparation training samples through a bi-directional generation process: (1) {Query-to-Pipeline Translation} to synthesize the core analytical logic, and (2) {Inverse Operator Injection} to simulate realistic data quality issues while guaranteeing recoverability.

\stitle{Task Instantiation.}
We first extract valid task tuples from the source datasets. Given a database and a ground-truth SQL query $q$, we identify the requisite source tables to form the clean input set $\mathcal{D}_{clean}$. Executing $q$ yields the target analytical table $T^*$.
To bridge the gap between SQL logic and natural language intent, we prompt an LLM to synthesize a {Target Schema Specification $\mathcal{S}$} based on the schema of $T^*$ and the semantic context of $q$.
This process yields a clean task tuple $(\mathcal{D}_{clean}, \mathcal{S}, T^*)$. The subsequent phases focus on synthesizing the missing operator pipeline and simulating realistic data entropy.

\subsection{Query-to-Pipeline Translation}
\label{subsec:forward_synthesis}

The objective of this phase is to construct a ground-truth operator pipeline $\Pi_{task}$ that transforms $\mathcal{D}_{clean}$ into $T^*$. Since existing datasets provide declarative SQL logic rather than imperative data manipulation steps, we must translate $q$ into an executable sequence within our operator space $\mathcal{O}$.

This translation is non-trivial. A single LLM often fails to decompose nested SQL queries involving multi-hop joins into correct atomic operators. To address this, we adopt a {voting-based multi-agent strategy} for Query-to-Pipeline translation.
Specifically, we instantiate $K$ distinct LLM agents (e.g., DeepSeek-R1, GPT-4o), each initialized with diverse reasoning prompts (e.g., ReAct~\cite{aksitov2023rest}, Plan-and-Solve~\cite{wang2024chain}), to generate a set of candidate pipelines $\mathbb{C} = \{\pi_1, \dots, \pi_K\}$.
We strictly verify each candidate $\pi_k$ by executing it on $\mathcal{D}_{clean}$. A pipeline is deemed valid if and only if its execution result $\hat{T}$ is identical to the ground-truth target $T^*$.
From the set of valid candidates, we select the pipeline with the minimum length as $\Pi_{task}$. This ensures that the synthesized training data represents the most efficient path to the solution, minimizing redundant operations.

\subsection{Inverse Operator Injection}
\label{subsec:backward_synthesis}

The forward generation phase yields a pipeline operating on pristine, structured data. However, real-world data lakes are notoriously ``dirty,'' plagued by missing values, formatting inconsistencies, and non-standard representations. Training an agent solely on clean data results in a policy that is brittle to noise.
To address this, we must simulate dirty source tables $\mathcal{D}_{dirty}$ from $\mathcal{D}_{clean}$.

Existing noise injection methods typically rely on rigid rules (e.g., random character deletion) or unconstrained LLM generation. The former lacks semantic realism, while the latter risks irreversible destruction (e.g., deleting a primary key), rendering the task unsolvable.
To this end, we propose {Inverse Operator Injection}, an iterative process that guarantees both \textit{diversity} and \textit{recoverability}.

\begin{algorithm}[!t]
\caption{\small Inverse Operator Injection Algorithm}
\label{alg:backward_synthesis}
\begin{algorithmic}[1]
    \Require \small Clean table $T_{clean}$, max iterations $L$, operator space $\mathcal{O}$
    \State $T_{dirty} \gets T_{clean},~\Pi_{clean} \gets []$
    \For{$i = 1$ to $L$}
        \State \commentstyle{Step 1: Select a cleaning operator as guidance}
        \State $\text{op}_{clean} \gets \mathtt{SampleOperator}(\mathcal{O}, T_{dirty})$
        
        \State \commentstyle{Step 2: Generate inverse logic (Dirtying)}
        \State $\text{code}_{inv} \gets \mathtt{GenerateInverseCode}(\text{op}_{clean})$
        \State $T' \gets \mathtt{ExecuteCode}(\text{code}_{inv}, T_{dirty})$
        
        \State \commentstyle{Step 3: Verify Recoverability}
        \State $\text{op}_{recov} \gets \mathtt{DeriveParams}(T', T_{dirty}, \text{op}_{clean})$
        \State $\hat{T} \gets \mathtt{Execute}(\text{op}_{recov}, T')$
        
        \If{$\mathtt{TableMatch}(\hat{T}, T_{dirty}) \land T' \neq T_{dirty}$}
            \State $T_{dirty} \gets T'$ \commentstyle{Commit dirty state}
            \State $\Pi_{clean} \gets [\text{op}_{recov}] + \Pi_{clean}$ \commentstyle{Prepend cleaning op}
        \EndIf
    \EndFor
    \State \Return $T_{dirty}$, $\Pi_{clean}$
\end{algorithmic}
\end{algorithm}

As outlined in Algorithm~\ref{alg:backward_synthesis}, we invert the standard data cleaning workflow. In each iteration, we first sample a cleaning operator $\text{op}_{clean}$ (e.g., \texttt{StandardizeDate}) from our operator space. We then prompt an LLM to generate a Python code snippet $\text{code}_{inv}$ that performs the \textit{inverse} operation—effectively injecting entropy. For instance, if $\text{op}_{clean}$ standardizes dates, $\text{code}_{inv}$ might convert ``2023-01-01'' into heterogeneous formats like ``01/01/23'' or ``Jan 1st, 2023''.

Crucially, we enforce a \textbf{Recoverability and Effectiveness Constraint} (Line 11). We verify that (1) the dirtying process actually altered the table ($T' \neq T_{dirty}$), and (2) we can derive specific parameters for $\text{op}_{clean}$ that can transform the dirty table $T'$ back to the state $T_{dirty}$. If and only if the execution of this recovery operator successfully restores the data, we accept $T'$ as the new state and prepend $\text{op}_{recov}$ to the cleaning pipeline $\Pi_{clean}$.

This process repeats for $L$ iterations. Finally, we concatenate the cleaning pipeline and the task pipeline to form the complete ground truth $\Pi_{total} = \Pi_{clean} \oplus \Pi_{task}$.
This synthesis engine produces high-quality training samples $(\mathcal{D}_{dirty}, \mathcal{S}, T^*, \Pi_{total})$ with pipeline lengths ranging from 1 to 28 operators. This dataset covers the full spectrum of DP tasks, providing the necessary complexity to train the agentic framework described in Section~\ref{sec:agentic_training_framework}.